% ----------------------------------
%   Appendix A: Implied Volatility Surface — Implementation Details
% ----------------------------------

\chapter{Implied Volatility Surface - Implementation Details}
\label{app:impl_vol}

This appendix describes the pipeline used to construct the SPY implied volatility
surface shown in Section~\ref{subsec:black_scholes_model},
including the practical difficulties encountered and the assumptions made at each step.
All code lives in \href{https://github.com/M315/TFM/tree/main/scripts/impl_vol/}{Implied Volatility Code} and is reproducible from the cached
market snapshot dated 2026-04-30.

% ----------------------------------
\section{Data}

SPY option chains are fetched via the \texttt{yfinance} library.
We restrict to maturities $4\,\text{wk} \leq T \leq 1\,\text{yr}$ and require a minimum
trading volume of 10 contracts per quote.
When both bid and ask are positive the mid-price is used;
otherwise the last traded price is taken as a fallback (relevant for pre-market snapshots
where the spread is stale but the last price is more informative).
Data for calls and puts are cached separately as \textit{parquet} files so that all
subsequent computations are fully reproducible without network access.

% ----------------------------------
\section{Risk-free rate}

The risk-free rate is taken from the US Treasury par yield curve published daily
by the US Department of the Treasury.
Par yields are available at standard maturities (1M, 2M, 3M, 4M, 6M, 1Y, 2Y, 3Y,
5Y, 7Y, 10Y, 20Y, 30Y) and are converted to continuously-compounded rates via
\begin{equation*}
    r_{\mathrm{cc}} = \ln\!\left(1 + \frac{y}{100}\right),
\end{equation*}
where $y$ is the par yield in percent.
Rates are linearly interpolated (in $T$) between the published nodes and held flat
outside the range.
Figure~\ref{fig:yield_curve} shows the resulting curve on the reference date.

\begin{figure}[h]
    \centering
    \includegraphics[width=0.72\textwidth]{yield_curve}
    \caption{US Treasury par yield curve on 2026-04-30 (log-maturity axis).
             Red dots are published nodes; the line is the linear interpolant
             used to assign a per-maturity discount rate.}
    \label{fig:yield_curve}
\end{figure}

Note that the risk-free rate could be refined further as the US Treasuries are not strictly risk-free,
as they carry a small credit risk premium.
A better discounting rate for SPY options would be the overnight index swap (OIS) rate rather than the Treasury curve,
but for our purposes the difference is not enough to justify the additional complexity of fetching and processing OIS data.

% ----------------------------------
\section{Dividend yield calibration}
\label{subsec:div_yield_cal}

SPY pays dividends quarterly, so any pricing formula must account for the dividend yield $q$.
Getting $q$ wrong breaks the internal consistency of the surface:
puts and calls at the same strike and maturity imply different volatilities,
producing a visible discontinuity at the money in any composite surface.
We went through three successive approximations before reaching a satisfactory result.

\paragraph{Stage 1 - $q = 0$.}
The simplest assumption is to ignore dividends entirely.
This is equivalent to treating SPY as a non-dividend-paying asset,
which is a standard simplification when the dividend yield is small relative to
the interest rate and the maturity is short.
For maturities beyond a few months, however, cumulative dividends are non-negligible
(SPY's trailing yield is approximately 1.1\%) and the forward price implied by
put-call parity,
\begin{equation*}
    F = S e^{(r - q)T},
\end{equation*}
is overestimated if $q$ is set to zero.
As a result, the Black-Scholes-Merton formula assigns different implied volatilities
to calls and puts at the same strike,
since each type is sensitive to the forward in a different direction.

\paragraph{Stage 2 - Fixed annual yield.}
The natural correction is to use the trailing twelve-month dividend yield reported
by the data provider ($q \approx 1.14\%$ for SPY on the reference date).
This constant is applied uniformly to all maturities.
The surface improves near the money, but the fix is imprecise:
SPY dividends are paid on a quarterly schedule, so the effective yield between
two option expirations depends on how many ex-dividend dates fall within that window.
A short-dated option expiring before the next ex-dividend date carries an effective
$q \approx 0$, while a longer-dated option spanning several dividend payments accumulates
a yield closer to the annual figure.
Applying a single annual constant systematically over-corrects for short maturities
and under-corrects in an uncontrolled way across the term structure.

\paragraph{Stage 3 - Put-Call parity calibration.}
The cleanest approach is to infer $q$ at each expiration directly from market prices
using put-call parity, keeping the Treasury rate $r$ fixed.
For each matched call-put pair at strike $K_i$ and expiration $T$,
put-call parity gives
\begin{equation*}
    C_i - P_i = S e^{-qT} - K_i e^{-rT}.
\end{equation*}
Defining the cost-of-carry forward $F = S e^{(r-q)T}$ and multiplying both sides
by $e^{rT}$ yields an implied forward per pair,
\begin{equation*}
    F_i = (C_i - P_i)\,e^{rT} + K_i,
\end{equation*}
where $q$ is encoded inside $F_i$ rather than appearing explicitly.
Taking the median over the relevant pairs at a given expiration yields a robust
forward estimate $\tilde{F}$, and inverting $\tilde{F} = S e^{(r-q)T}$ gives
\begin{equation*}
    q = r - \frac{\ln(\tilde{F}/S)}{T}.
\end{equation*}

A key implementation choice is \emph{which pairs to include}.
Put-call parity holds model-free for any European options,
but the quality of the forward estimate depends on how cleanly each leg is priced.
For a strike $K < S$ the call is in-the-money and the put is out-of-the-money;
for $K > S$ the roles are reversed.
Deep in-the-money options are significantly less liquid:
their price is dominated by intrinsic value, the bid-ask spread is wide relative
to the time value, and last-trade prices can be hours stale.
Feeding these pairs into the forward estimate biases $\tilde{F}$ toward noisy values.

We therefore restrict the calibration to near-ATM pairs satisfying
$|K/S - 1| \leq 0.05$ (i.e.\ strikes within $\pm 5\%$ of the current spot).
In this band, both the call and the put are close to at-the-money;
both legs have tight spreads and meaningful open interest,
and the time value is large relative to the bid-ask noise.
The median over these pairs (rather than the mean or a least-squares fit)
guards against the occasional stale quote even in the near-ATM region.
If fewer than two near-ATM pairs are available the constraint is relaxed to all matched pairs.

Since $r$ is fixed from the Treasury curve, there is only one free parameter per
expiration, which avoids the instability of simultaneously fitting both $r$ and $q$.
Calibrated values are bounded to $q \in [-0.02, 0.10]$; if the result is outside
this range (a sign of too few matched pairs or erroneous data), the provider value
is used as a fallback.

The resulting term structure shows $q \approx 0$ for near-term expirations
(the next quarterly dividend has not yet been declared within the window)
and a gradual convergence to the trailing annual yield for longer maturities,
which is economically consistent with the quarterly payment schedule.

Figure~\ref{fig:atm_gap} illustrates the three stages for the December 2026 expiration
($T \approx 0.63$ yr), where the effect of the dividend yield is large enough to be
clearly visible.
Each panel overlays the put-implied and call-implied volatility curves across the
full strike range.
With $q = 0$ (Stage 1) the two curves diverge significantly and the level is
shifted upward for OTM puts and downward for OTM calls.
With a fixed annual $q$ (Stage 2) the ATM region improves but the term-structure
mismatch remains.
With the PCP-calibrated $q$ (Stage 3) the two curves sit on top of each other
near the money, confirming that the composite surface can be joined smoothly.
The residual divergence at deep in-the-money strikes visible in all three panels
is due to liquidity issues (wide bid-ask spreads, stale prices) rather than the
dividend yield assumption, and is removed by restricting the composite to
out-of-the-money quotes only.

\begin{figure}[h]
    \centering
    \includegraphics[width=\textwidth]{annex_atm_gap}
    \caption{Call-implied (solid blue) and put-implied (dashed red) volatility across
             all available strikes at the December 2026 expiration ($T \approx 0.63$ yr),
             under three dividend yield assumptions.
             Stage 1 ($q=0$): curves diverge sharply — the forward is overestimated,
             shifting call and put IVs in opposite directions.
             Stage 2 (fixed annual $q$): improved near ATM, but the flat term structure
             leaves systematic errors for maturities where the true $q$ differs from
             the annual average.
             Stage 3 (PCP-calibrated $q$ from near-ATM pairs): the two curves coincide
             in the OTM region, where both legs are liquid.
             The residual divergence far from ATM (deep ITM options) is a liquidity
             artefact — those strikes are excluded from the composite surface.}
    \label{fig:atm_gap}
\end{figure}

% ----------------------------------
\section{OTM composite construction}

Once $r$ and $q$ are consistent with market prices, the composite surface is formed
by taking out-of-the-money options only:
OTM puts ($K < S$) for the left wing and OTM calls ($K \geq S$) for the right wing.

The rationale is that in-the-money options are significantly less liquid than
their out-of-the-money counterparts: ITM bids tend to be stale, spreads are wide,
and the option price is dominated by intrinsic value, so a small absolute pricing
error translates into a large relative error in the time value — and therefore a
large error in the inverted implied volatility.
This is visible in Figure~\ref{fig:atm_gap}: the divergence between call and put
IV curves away from the money persists even in Stage~3, not because $q$ is wrong
but because deep-ITM prices are unreliable.

Figures~\ref{fig:iv_calls} and \ref{fig:iv_puts} show the full call and put surfaces
(all strikes, no ITM filter) to illustrate the noise.
The clean composite (Figure~\ref{fig:iv_surface}) is obtained after restricting to
$0.75 \leq K/S \leq 1.25$ and keeping only the OTM leg for each option type.

\begin{figure}[h]
    \centering
    \begin{minipage}{0.49\textwidth}
        \centering
        \includegraphics[width=\textwidth]{iv_calls}
        \caption{IV surface from all call quotes (including deep ITM, $K \ll S$).
                 Deep ITM calls carry noisy implied volatilities due to wide spreads
                 and stale prices.}
        \label{fig:iv_calls}
    \end{minipage}
    \hfill
    \begin{minipage}{0.49\textwidth}
        \centering
        \includegraphics[width=\textwidth]{iv_puts}
        \caption{IV surface from all put quotes (including deep ITM, $K \gg S$).
                 Same liquidity issue for strikes far above the spot.}
        \label{fig:iv_puts}
    \end{minipage}
\end{figure}

% ----------------------------------
\section{Expiration selection}

SPY has a dense cluster of weekly expirations in the short end (roughly 1-8 weeks).
Including all available expirations would over-weight this region and produce a surface
that appears smooth in absolute maturity but poorly samples the longer end.
To obtain a balanced representation, the full maturity range is divided into
$N = 10$ log-spaced bins; within each bin the expiration with the highest combined
open interest (calls plus puts) is kept.
This ensures that liquidity guides the selection and that the resulting sample
is approximately uniform in $\log T$, which is the natural scale on which the
term structure of volatility varies.

% ----------------------------------
\section{Summary of assumptions}

For reference, the key modelling assumptions are:
\begin{itemize}
    \item \textbf{European exercise}: SPY options are technically American-style.
          The early-exercise premium is small for an index ETF with a low continuous yield
          and is neglected here; the BSM formula is used as an approximation.
    \item \textbf{Continuous dividends}: actual SPY dividends are discrete quarterly
          payments. The PCP calibration implicitly converts them into a continuous yield
          that is consistent with the current option prices, which is the standard
          market convention.
    \item \textbf{Moneyness filter} $0.75 \leq K/S \leq 1.25$: strikes outside
          this range are too illiquid to contribute reliable IV observations.
    \item \textbf{Volume filter}: quotes with fewer than 10 contracts of daily volume
          are excluded.
    \item \textbf{Single snapshot}: the surface is computed from a single end-of-day
          snapshot; intraday bid-ask movements are not captured.
\end{itemize}
